\documentclass[a4paper,12pt]{article}
\usepackage[T2A]{fontenc}
\usepackage[utf8]{inputenc}
\usepackage[russian]{babel}
\usepackage{amsmath,amssymb,mathtools}
\usepackage{helvet}
\renewcommand{\familydefault}{\sfdefault}
\usepackage{icomma}
\usepackage[a4paper,left=22mm,right=22mm,top=20mm,bottom=22mm]{geometry}
\usepackage{microtype}
\usepackage{tikz}
\usepackage{tabularx,booktabs,longtable}
\usepackage{enumitem}
\usepackage{xcolor}
\usepackage{hyperref}
\usepackage{parskip}
\usepackage{fancyhdr}
\usepackage{setspace}
\definecolor{accent}{HTML}{7D2835}
\definecolor{darkaccent}{HTML}{4D1820}
\definecolor{textgray}{HTML}{333333}
\definecolor{softgray}{HTML}{EEE8DF}
\hypersetup{colorlinks=true,linkcolor=darkaccent,urlcolor=darkaccent,pdfauthor={Лёвин Артём Александрович},pdftitle={Кредиты: аннуитетная и дифференцированная схемы}}
\onehalfspacing
\setlength{\parindent}{0pt}
\setlength{\headheight}{25pt}
\renewcommand{\arraystretch}{1.3}
\emergencystretch=2em
\pagestyle{fancy}
\fancyhf{}
\fancyhead[L]{\small\textcolor{textgray}{Финансовая математика}}
\fancyhead[R]{\small\textcolor{textgray}{Николь Саркисьянц}}
\fancyfoot[C]{\small\thepage}
\renewcommand{\headrulewidth}{0.4pt}
\newcommand{\sectiontitle}[1]{\vspace{0.4em}{\Large\bfseries\textcolor{darkaccent}{#1}}\par\vspace{0.2em}{\color{accent}\rule{\linewidth}{0.7pt}}\par\vspace{0.45em}}
\newcommand{\important}[1]{\textbf{\textcolor{darkaccent}{#1}}}
\begin{document}
\begin{titlepage}\thispagestyle{empty}
\begin{center}\vspace*{2.2cm}
{\Large\textcolor{accent}{\bfseries ЧЕК-ЛИСТ}}\\[0.8cm]
{\huge\bfseries\textcolor{darkaccent}{Кредиты}}\\[0.3cm]
{\LARGE\bfseries Аннуитетная и дифференцированная схемы}\\[1.5cm]
{\large Подготовка к ЕГЭ: модель, формулы, алгоритмы и практика}\\[1.8cm]
{\large\bfseries Лёвин Артём Александрович}\\[0.25cm]
{\large эксклюзивно для Николь Саркисьянц}\\[1.2cm]
{\large 28 июля 2026 года}
\vfill
\begin{tikzpicture}[remember picture,overlay]
\draw[line width=1.2pt,color=accent] ([xshift=16mm,yshift=20mm]current page.south west)
.. controls ([xshift=58mm,yshift=34mm]current page.south west) and ([xshift=-65mm,yshift=7mm]current page.south east) .. ([xshift=-16mm,yshift=20mm]current page.south east);
\end{tikzpicture}
\end{center}\end{titlepage}

\sectiontitle{1. Универсальная модель кредита}
Обозначения:
\[
S\text{ --- первоначальная сумма},\qquad p=\frac{P}{100},\qquad r=1+p,\qquad n\text{ --- число периодов}.
\]
В каждом периоде выполняется одна и та же цепочка:
\[
\text{долг}\longrightarrow\text{начисление процентов}\longrightarrow\text{общая задолженность}\longrightarrow\text{платёж}\longrightarrow\text{остаток}.
\]
После последнего платежа остаток задолженности равен нулю.

\sectiontitle{2. Аннуитетная схема}
Платёж $x$ одинаков в каждом периоде. Если $D_0=S$, то
\[
\boxed{D_k=rD_{k-1}-x}.
\]
Первые остатки:
\[
D_1=Sr-x,\qquad D_2=Sr^2-x(r+1),\qquad D_3=Sr^3-x(r^2+r+1).
\]
Общая формула:
\[
D_k=Sr^k-x\bigl(r^{k-1}+r^{k-2}+\dots+r+1\bigr).
\]
Из условия $D_n=0$ получаем
\[
\boxed{x=\frac{Sr^n}{r^{n-1}+r^{n-2}+\dots+r+1}}.
\]
При $r\ne1$ можно использовать сумму геометрической прогрессии:
\[
x=\frac{Sr^n(r-1)}{r^n-1}.
\]
При $r=1$ делить на $r-1$ нельзя; тогда $D_k=S-kx$.

\sectiontitle{3. Выплаты и переплата}
Для аннуитетной схемы
\[
\text{общая сумма выплат}=nx,
\qquad
\text{переплата}=nx-S,
\]
\[
\boxed{\text{процент переплаты}=\frac{nx-S}{S}\cdot100\%}.
\]
При фиксированных $r$ и $n$ процент переплаты не зависит от $S$: множитель $S$ сокращается.

\sectiontitle{4. Пример аннуитетной схемы}
Два платежа по $121\,000$ рублей, ставка $10\%$. Тогда $r=1{,}1$ и
\[
Sr^2-x(r+1)=0,
\qquad
S=\frac{x(r+1)}{r^2}=\frac{121\,000\cdot2{,}1}{1{,}21}=210\,000.
\]
Общая сумма выплат $242\,000$ рублей, переплата $32\,000$ рублей.

\sectiontitle{5. Дифференцированная схема}
Тело кредита погашается равными долями:
\[
\boxed{\frac{S}{n}\text{ в каждом периоде}}.
\]
Перед $k$-м платежом основной долг равен
\[
B_{k-1}=S\frac{n-k+1}{n}.
\]
Следовательно,
\[
\boxed{A_k=\frac{S}{n}+pB_{k-1}=\frac{S}{n}\bigl(1+p(n-k+1)\bigr)}.
\]
Платежи образуют убывающую арифметическую прогрессию:
\[
A_1>A_2>\dots>A_n,
\qquad
A_k-A_{k+1}=\frac{pS}{n}.
\]
Общая сумма выплат и переплата:
\[
\boxed{\text{общая сумма}=S+\frac{pS(n+1)}{2}},
\qquad
\boxed{\text{переплата}=\frac{pS(n+1)}{2}}.
\]

\sectiontitle{6. Сравнение схем}
\begin{table}[h]\centering
\begin{tabularx}{\linewidth}{@{}lXX@{}}\toprule
Признак & Аннуитетная & Дифференцированная\\\midrule
Что постоянно & полный платёж $x$ & доля тела кредита $S/n$\\
Платежи & одинаковые & уменьшаются\\
Первый платёж & равен остальным & самый большой\\
Главная модель & $D_k=rD_{k-1}-x$ & $A_k=S/n+pB_{k-1}$\\\bottomrule
\end{tabularx}\end{table}

\sectiontitle{7. Алгоритм решения}
\begin{enumerate}[leftmargin=2em,itemsep=0.3em]
\item Выпишите $S$, $P$, $p$, $r$, $n$ и искомую величину.
\item Определите схему по формулировке условия.
\item Составьте таблицу изменения задолженности.
\item Сначала выведите формулу в буквенном виде.
\item Используйте условие полного погашения.
\item Подставьте числа и выполните расчёты.
\item Найдите общую сумму выплат, переплату и процент переплаты.
\item Проверьте экономический смысл ответа.
\end{enumerate}

\sectiontitle{8. Типичные ошибки}
\begin{itemize}[leftmargin=1.7em,itemsep=0.3em]
\item Путать $P\%$ и $p=P/100$.
\item Забывать, что $r=1+p$.
\item Подставлять числа до построения модели.
\item В дифференцированной схеме брать долю от текущего остатка вместо $S/n$.
\item Не вычитать $S$ при вычислении переплаты.
\item Не проверять условие $D_n=0$ и смысл найденных корней.
\end{itemize}

\clearpage
\sectiontitle{Тренировочный блок}
\begin{enumerate}[leftmargin=2em,itemsep=0.75em]
\item Кредит $500\,000$ рублей выдан под $8\%$. Запишите $r$, $D_1$, $D_2$, $D_3$ для равного платежа $x$.
\item Два равных платежа по $121\,000$ рублей, ставка $10\%$. Найдите $S$, общую сумму выплат и переплату.
\item Кредит $120\,000$ рублей погашается за $4$ года равными долями под $20\%$. Найдите платежи.
\item Кредит $300\,000$ рублей погашается тремя равными платежами под $10\%$. Найдите $x$ и переплату.
\item Три равных платежа по $133\,100$ рублей, ставка $10\%$. Найдите $S$.
\item Кредит $210\,000$ рублей погашается двумя платежами по $121\,000$ рублей. Найдите ставку.
\item Кредит $500\,000$ рублей погашается за $5$ лет равными долями под $12\%$. Найдите платежи и общую сумму.
\item В дифференцированной схеме на $5$ лет первый платёж $150\,000$, последний $110\,000$. Найдите $S$ и ставку.
\item Кредит $400\,000$ рублей выдан на $4$ года под $10\%$. Сравните общие суммы выплат по двум схемам.
\item Докажите, что при одинаковых $r$ и $n$ процент переплаты по аннуитетной схеме не зависит от $S$.
\end{enumerate}

\clearpage
\sectiontitle{Ответы}
\begin{longtable}{@{}p{0.07\linewidth}p{0.87\linewidth}@{}}\toprule
№&Ответ\\\midrule\endfirsthead\toprule №&Ответ\\\midrule\endhead\bottomrule\endfoot
1&$r=1{,}08$; $D_1=540\,000-x$; $D_2=583\,200-2{,}08x$; $D_3=629\,856-3{,}2464x$.\\
2&$S=210\,000$ руб.; общая сумма $242\,000$ руб.; переплата $32\,000$ руб.\\
3&$54\,000$, $48\,000$, $42\,000$, $36\,000$ руб.\\
4&$x\approx120\,634{,}44$ руб.; переплата $\approx61\,903{,}32$ руб.\\
5&$S=331\,000$ руб.\\
6&$10\%$.\\
7&$160\,000$, $148\,000$, $136\,000$, $124\,000$, $112\,000$ руб.; общая сумма $680\,000$ руб.\\
8&$S=500\,000$ руб.; ставка $10\%$.\\
9&Аннуитетная сумма $\approx504\,753{,}29$ руб.; дифференцированная $500\,000$ руб.; разница $\approx4\,753{,}29$ руб.\\
10&После подстановки $x$ множитель $S$ сокращается; остаётся выражение, зависящее только от $r$ и $n$.\\
\end{longtable}
\end{document}
